\section{Preliminaries}
\label{sec:prelim}

We briefly recap the setup of \citet{boyarsky2026}. The hierarchical
data-generating process is
\begin{equation}
\label{eq:dgp}
Y \;=\; A\,\tau(X) + C_m^\top\beta_1 + C_r^\top\beta_2 + g_0(X)
 + C_m^\top\gamma_1 + C_r^\top\gamma_2 + \varepsilon,
\end{equation}
where $A\in\{0,1\}$ is the treatment, $X\in\mathbb{R}^{d_X}$ individual-level
covariates, $C_m\in\mathbb{R}^{d_M}$ site-level main effects, and
$C_r\in\mathbb{R}^{d_I}$ site-level interaction covariates. The parameter
$\theta := (\beta_1,\beta_2,\gamma_1,\gamma_2)\in\mathbb{R}^{d_g}$ with
$d_g = 2(d_M + d_I)$ captures direct and indirect effects.

\paragraph{Proposition 1 of \protect\citet{boyarsky2026}.}
\emph{Under a constant-propensity RCT and mean-zero errors, for a target profile
$(t_m,t_r)$ with $\mathrm{Pr}[C_m=t_m, C_r=t_r]>0$,}
\[
\psi(\theta;\gamma) \;:=\; \mathbb{E}[Y(1)-Y(0)\,|\,C_m=t_m, C_r=t_r]
\]
\emph{admits the representation}
\begin{equation}
\label{eq:rep}
\psi(\theta;\gamma)
 \;=\; c_0(\gamma) + \beta_1^\top \tilde c_m(\gamma) + \beta_2^\top \tilde c_r(\gamma),
\qquad
\Sigma\theta + b = 0,
\end{equation}
\emph{with the moment matrix $\Sigma$ of Equation~(A.9) in their appendix, namely a
block matrix of $\mathbb{E}_X\!\bigl[\mathrm{Cov}(C_\bullet,C_\bullet \mid X)\bigr]$
entries with $\bullet \in \{m, r\}$.}

The key structural fact we exploit in this paper is that $\Sigma$ is a linear
functional of the joint distribution of $(C_m, C_r, X)$, and in particular its
rank is governed by the rank of the $(S \times (d_M + d_I))$ matrix whose rows
are the site-level contexts $(C_m^{(s)}, C_r^{(s)})$, weighted by site sample
sizes and individual covariate overlap. When $S$ is small relative to $d_M +
d_I$, this matrix, and hence $\Sigma$, is rank-deficient. This is precisely why
the \citet{meager2019understanding} micro-credit application with $S=7$ sites
has non-trivial partial identification.

\paragraph{Partial-identification bounds.}
\citet{boyarsky2026} define
\begin{equation}
\label{eq:bounds}
v_{\ell}(\gamma) \;=\; \min_\theta\,\bigl\{\psi(\theta;\gamma)\,:\,
 \Sigma\theta + b = 0,\; r(\theta)\le 0\bigr\},
\qquad
v_{u}(\gamma) \;=\; \max_\theta\,\bigl\{\psi(\theta;\gamma)\,:\,
 \Sigma\theta + b = 0,\; r(\theta)\le 0\bigr\},
\end{equation}
and the width of the partial-identification interval at target $(t_m,t_r)$ is
\begin{equation}
\label{eq:width}
w_r(t_m, t_r; \Sigma, b) \;=\; v_u(\gamma) - v_\ell(\gamma).
\end{equation}
We suppress the dependence of $c_0, \tilde c_m, \tilde c_r$ on $\gamma$ when no
confusion arises. None of these interact with the feasible set, but they play
different roles in the width: the additive constant $c_0(\gamma)$ shifts both
endpoints equally and therefore cancels in \eqref{eq:width}, whereas the linear
target direction $c(t)$ (built from $\tilde c_m, \tilde c_r$ and the target
profile) determines both the endpoints and the width.

\paragraph{Restrictions.}
The analyst chooses $r(\theta)\le 0$ from a library of affinely-representable
restrictions (effect-ordering, sign constraints, norm bounds, monotonicity).
Throughout this paper we assume $r$ defines a compact polyhedron $\mathcal{R} :=
\{\theta : r(\theta)\le 0\}$. This is the regime where the width is finite and
the LP/MIP program is well-posed. The extension to relaxed
restriction sets (where unrestricted widths are permitted) is discussed in
Section~\ref{sec:extensions}.
